% 連続体力学 連続体力学-4 問題3（029-hooke） 解いた日 2026-11-25
%% kotae
\toi{1}{◎}{}
\begin{align*}
P_{ij}&=\sum_{k,\ell}C_{ijk\ell}E_{k\ell}\\
&=\sum_{k,\ell}\bigl(\lambda\delta_{ij}\delta_{k\ell}\\
&\qquad+\mu(\delta_{ik}\delta_{j\ell}+\delta_{i\ell}\delta_{jk})\bigr)E_{k\ell}\\
&=\sum_{k}\lambda E_{kk}\delta_{ij}+\underbrace{\mu(E_{ij}+E_{ji})}_{\text{$\tens{E}$ は対称テンソル}}\\
&=\lambda\Tr(\tens{E})\delta_{ij}+2\mu E_{ij}
\end{align*}
\toi{2}{◎}{}
\[ \lambda\Tr(\tens{E})=60\times0.2=\kakimachigai{\qty{12}{\giga\pascal}}{\qty{12}{\mega\pascal}} \]
\[ P_{ij}= \]
\[ \setlength{\arraycolsep}{2pt}\begin{pmatrix}50\times1.0+12&50\times0.5&0\\50\times0.5&50\times(-0.4)+12&0\\0&0&50\times(-0.4)+12\end{pmatrix} \]
\begin{align*}
&\qquad\times10^{-3}\,\unit{\giga\pascal}\\
&=\begin{pmatrix}62&25&0\\25&-8.0&0\\0&0&-8.0\end{pmatrix}\unit{\mega\pascal}
\end{align*}
\toi{3}{○}{単位}
\begin{align*}
\Tr(\tens{P})&=\kakimachigai{\sum_{i}}{\sum_{i,j}}P_{ij}\delta_{ij}\\
&=\sum_{i}(\lambda\Tr(\tens{E})\delta_{ij}+2\mu E_{ij}\delta_{ij})\\
&=3\lambda\Tr(\tens{E})+2\mu\Tr(\tens{E})\\
&=(3\lambda+2\mu)\Tr(\tens{E})
\end{align*}
\[ \Tr(\tens{P})=62-8-8=46\nuke{\unit{\mega\pascal}} \]
\begin{align*}
&(3\lambda+2\mu)\times0.2\\
&=(180+50)\times\frac{1}{5}\\
&=46\nuke{\unit{\mega\pascal}}
\end{align*}
\naze{数値には単位が要る。両辺とも $\qty{46}{\mega\pascal}$ で，右辺は $\qty{230}{\giga\pascal}\times\num{0.20e-3}$ である（$0.2$ は $\num{0.20e-3}$ の $10^{-3}$ を省いた形）。式の変形と数値の一致は正しい。}
%% kizuki
\item 歪みの $10^{-3}$ を頭の中で持ち運ぶと，(2) の $\lambda\Tr(\tens{E})=\qty{12}{\giga\pascal}$ や (3) の単位なしの 46 のように，書いた数値と単位がずれる。$\qty{1}{\giga\pascal}\times10^{-3}=\qty{1}{\mega\pascal}$ と書いてから計算し，数値の答えには毎回単位を付ける。
